% !TEX root = ../main.tex
\chapter{Introduzione}\label{chap:introduction}

Le architetture a microservizi cloud-native hanno reso le \gls{API} REST il principale confine operativo tra sistemi: ogni servizio espone un insieme di endpoint con le proprie regole di autenticazione, autorizzazione e validazione dei dati, e la superficie di attacco complessiva cresce proporzionalmente alla proliferazione di questi endpoint. Valutarne la sicurezza richiede strumenti che conoscano il contratto che ogni richiesta dovrebbe rispettare; gli scanner tradizionali, privi di quella conoscenza, operano sulla sintassi e rimangono ciechi alla classe di vulnerabilità semantiche che caratterizza le \gls{API} REST. Questo lavoro progetta, implementa e valida empiricamente un framework di security assessment che usa la specifica OpenAPI come oracolo formale per costruire verifiche riproducibili e prioritizzate su qualsiasi sistema REST documentato, senza dipendenze da un target specifico.


% ─────────────────────────────────────────────────────────────────────────────
\section{Motivazione e Contesto}\label{sec:motivazione}
% ─────────────────────────────────────────────────────────────────────────────

L'adozione delle architetture a microservizi cloud-native ha redistribuito la complessità applicativa su un numero crescente di servizi autonomi, ciascuno responsabile di un sottodominio del sistema e comunicante con gli altri tramite \gls{API} REST. La conseguenza diretta è la moltiplicazione della superficie esposta: ogni servizio aggiunge endpoint al perimetro, spesso in assenza di un inventario centralizzato e di una governance sistematica del loro ciclo di vita. In questo contesto il gateway \gls{API} è emerso come punto di enforcement centralizzato: posto tra i client e i servizi interni, concentra in un unico piano di configurazione ispezionabile politiche di autenticazione, autorizzazione, rate limiting e altre funzioni di controllo, tra le molte che un deployment reale può richiedere. Verificare se il gateway applica correttamente le politiche dichiarate e se gli endpoint esposti corrispondono a quelli documentati è parte dell'obiettivo di un sistema di security assessment; ma le vulnerabilità rilevanti non si limitano alla configurazione del piano di controllo. Molte emergono interrogando la logica semantica degli endpoint stessi, in modo indipendente dal gateway che li espone.

L'\gls{OWASP} API Security Top~10~\cite{MISSING:owasp:api-security:2023} % TODO: add to Zotero
formalizza le dieci categorie di rischio più frequenti nelle \gls{API} REST, con le prime posizioni occupate da vulnerabilità semantiche, tra cui violazioni di autorizzazione a livello di oggetto (\gls{BOLA}), difetti di autenticazione ed esposizione non intenzionale di proprietà. Queste vulnerabilità si manifestano come richieste formalmente corrette che violano le regole che il sistema impone su chi può accedere a cosa: quali risorse un client autenticato può leggere o modificare, con quali credenziali, e in quale sequenza di operazioni; per questa ragione non producono anomalie sintattiche né nella forma della richiesta né nella struttura della risposta, e sfuggono agli strumenti che cercano esclusivamente questi segnali. Le \gls{API} hanno consolidato il proprio ruolo di vettore primario negli incidenti di sicurezza enterprise~\cite{MISSING:salt-security:2024}: % TODO: add to Zotero
la frequenza di violazioni correlata a configurazioni errate dei gateway e a controlli di autorizzazione insufficienti segnala un disallineamento strutturale tra la velocità con cui le \gls{API} vengono esposte e la maturità degli strumenti disponibili per valutarne la sicurezza.

Questa complessità semantica rende difficile affiancare al ciclo di sviluppo continuativo una valutazione sistematica della sicurezza. Un penetration test manuale produce un'osservazione approfondita ma circoscritta al momento e al contesto in cui viene condotto; non si presta a essere ripetuto automaticamente a ogni variazione del target, né a produrre risultati confrontabili tra esecuzioni successive senza intervento umano. Dotare il processo di assessment di riproducibilità deterministica, tracciabilità dei verdetti agli standard normativi e integrabilità nelle pipeline di sviluppo continuativo richiede che le verifiche siano parametrizzate sul contratto del target e indipendenti dall'operatore che le esegue; solo a quella condizione possono essere eseguite in forma autonoma su qualsiasi sistema documentato.


% ─────────────────────────────────────────────────────────────────────────────
\section{Definizione del Problema e Gap nella Ricerca}\label{sec:gap-ricerca}
% ─────────────────────────────────────────────────────────────────────────────

Tra le limitazioni degli approcci esistenti al security assessment delle \gls{API} REST, tre si sono rivelate centrali per questo lavoro: ciascuna è radicata in una scelta di design diversa e porta implicazioni architetturali distinte.

\paragraph{G\textsubscript{1} Cecità Semantica}\label{g1:cecita-semantica}
Gli scanner \gls{DAST} e i fuzzer generici operano sulla sintassi delle richieste \gls{HTTP}: generano payload, osservano i codici di risposta e cercano anomalie nella forma della request (come le injection sui parametri) o nella struttura della response (come codici di errore inattesi). Le vulnerabilità semantiche delle \gls{API} REST si manifestano invece come richieste sintatticamente corrette che violano le regole di accesso del sistema: quali risorse un client può richiedere, con quali credenziali, rispetto a quali altri soggetti. Per rilevarle è necessario conoscere il contratto che regola ogni interazione, ad esempio quali ruoli possono accedere a quale risorsa, quale sia il ciclo di vita atteso di un token, e quali sequenze di operazioni abbiano senso semantico. Uno scanner privo di quella conoscenza non può costruire sistematicamente le probe che raggiungono questa logica.

\paragraph{G\textsubscript{2} Portabilità e Profondità Semantica}\label{g2:portabilita-profondita}
Gli strumenti specializzati di security assessment per \gls{API} presentano un compromesso ricorrente: quelli progettati per un gateway o una piattaforma specifica offrono verifiche approfondite ma sono applicabili solo a quel contesto; quelli sufficientemente generici da operare su qualsiasi target, usano la \gls{OAS} principalmente per generare richieste valide dal punto di vista sintattico, senza impiegarla come oracolo per valutare se le risposte rispettino il contratto di sicurezza dichiarato. Il risultato è che profondità e portabilità vengono ottimizzate separatamente, e la specifica del target costituisce raramente l'oracolo condiviso tra le due esigenze. Questo riflette anche un limite dell'ecosistema: l'adozione di specifiche formali è ancora disomogenea nel panorama reale, e quando esistono non sempre riflettono fedelmente l'implementazione effettiva, rendendo difficile affidarsi ad esse come fonte autoritativa per la valutazione della sicurezza.

\paragraph{G\textsubscript{3} Riproducibilità Deterministica}\label{g3:riproducibilita}
Un assessment condotto da un operatore in un momento specifico produce un'osservazione il cui esito dipende tanto dalle caratteristiche del target quanto dalle scelte metodologiche di chi lo ha condotto. In un processo di sviluppo continuativo, dove il sistema evolve costantemente e ogni modifica può introdurre regressioni di sicurezza, questa dipendenza dall'operatore impedisce di affiancare al ciclo di rilascio una verifica automatica e confrontabile nel tempo: un operatore umano non può seguire la frequenza dei cambiamenti, e la mancanza di determinismo nell'output rende impossibile distinguere una regressione reale da una variazione metodologica. Senza un meccanismo che renda il comportamento del tool indipendente dal contesto di esecuzione e verificabile su run successive, la riproducibilità rimane un'assunzione implicita piuttosto che una garanzia dimostrabile.


% ─────────────────────────────────────────────────────────────────────────────
\section{Obiettivi della Ricerca}\label{sec:obiettivi}
% ─────────────────────────────────────────────────────────────────────────────

Ai tre gap identificati corrispondono tre obiettivi di ricerca. Ciascuno è formulato attorno a una proprietà concreta che le scelte architetturali realizzano e che la validazione empirica può valutare.

\paragraph{O\textsubscript{1} Tool Contract-Driven Agnostico}\label{o1:framework-agnostico}
Il primo obiettivo, in risposta al \nameref{g1:cecita-semantica}, è progettare e implementare un tool che adotti il paradigma contract-driven: la \gls{OAS} del target funge da oracolo formale per costruire verifiche semanticamente significative, mentre il file di configurazione fornisce la conoscenza concreta del deployment specifico, come i valori degli endpoint parametrici e le credenziali dei ruoli. La combinazione dei due mira a eliminare la necessità di riferimenti hardcoded a qualsiasi sistema specifico e a organizzare le verifiche in una tassonomia strutturata di domini di sicurezza con priorità esplicite; ogni test sarà corredato dai relativi riferimenti normativi, producendo evidenza tracciabile verso categorie \gls{OWASP}, codici \gls{CWE} e standard applicabili.

\paragraph{O\textsubscript{2} Superamento del Compromesso Portabilità-Profondità}\label{o2:profondita-portabilita}
Il secondo obiettivo, in risposta al \nameref{g2:portabilita-profondita}, è una conseguenza diretta dell'approccio adottato in \nameref{o1:framework-agnostico}: un tool che deriva tutta la conoscenza del target dalla specifica e dalla configurazione, senza dipendenze hardcoded, è per costruzione applicabile a qualsiasi sistema REST documentato senza modifiche al codice. La profondità semantica delle verifiche non dipende quindi dalla piattaforma su cui il target è deployato, ma dalla qualità della specifica che lo descrive.

\paragraph{O\textsubscript{3} Determinismo e Integrazione nel Ciclo Continuativo}\label{o3:riproducibilita-cicd}
Il terzo obiettivo, in risposta al \nameref{g3:riproducibilita}, è progettare il tool in modo che esecuzioni successive con la stessa configurazione sullo stesso target producano output identici, indipendentemente dall'operatore, e verificare questa proprietà empiricamente. Il determinismo dell'output è la condizione che rende la verifica affiancabile al ciclo di sviluppo continuativo tramite exit code semantici, senza richiedere interpretazione manuale dei risultati tra un'esecuzione e la successiva.


% ─────────────────────────────────────────────────────────────────────────────
\section{Organizzazione della Tesi}\label{sec:organizzazione}
% ─────────────────────────────────────────────────────────────────────────────

La tesi è organizzata nei seguenti capitoli.

Nel \textbf{Capitolo~\ref{chap:background}} viene costruito il contesto tecnico che i capitoli successivi presuppongono. Si analizzano l'evoluzione verso le architetture a microservizi e i principali stili di comunicazione con le loro implicazioni di sicurezza; il ruolo del gateway \gls{API} come punto di enforcement centralizzato e i criteri che ne motivano la scelta; il divario strutturale tra gli strumenti di testing tradizionali e le vulnerabilità semantiche delle \gls{API} REST; e l'evoluzione degli approcci di testing dal fuzzing cieco agli strumenti contract-driven.

Nel \textbf{Capitolo~\ref{chap:methodology}} vengono presentati la metodologia e l'architettura del tool, con un approccio top-down che parte dal livello di sistema per scendere ai componenti. Si introducono agnosticismo applicativo, paradigma contract-driven e riproducibilità deterministica come fondamenta del design; si descrive l'architettura dell'assessment engine con il suo scheduling basato su \gls{DAG}, e si definisce la tassonomia degli otto domini di sicurezza con il box gradient come mapping tra precondizioni di accesso e profondità dell'assessment.

Nel \textbf{Capitolo~\ref{chap:implementation}} le proprietà architetturali del capitolo precedente vengono portate al livello implementativo. Si descrive la pipeline di esecuzione a sette fasi, il meccanismo di discovery dinamico dei test, il contratto delle classi base \texttt{BaseTest} ed \texttt{ExternalToolTest}, la gerarchia dei connector per l'integrazione di tool esterni, e i meccanismi di gestione sicura delle credenziali e di streaming dell'evidence store.

Nel \textbf{Capitolo~\ref{chap:experiments}} viene condotta la validazione sperimentale del sistema su un target reale: Forgejo~14.0.3 esposto tramite Kong~3.9 in modalità DB-less. La validazione si articola su quattro livelli: la composizione del testbed, la qualità ingegneristica del codebase come prerequisito di affidabilità dello strumento di misura, la copertura funzionale e l'idempotenza su due run indipendenti, e infine il profilo computazionale dell'assessment.

Nel \textbf{Capitolo~\ref{chap:discussion}} i limiti strutturali del paradigma contract-driven vengono ricondotti alle scelte architetturali che li originano; si analizzano le condizioni concrete di integrabilità nelle pipeline \gls{CI/CD} industriali e la variabilità di copertura nei diversi contesti di deployment; e infine si tracciano le traiettorie di sviluppo futuro distinte in estensioni operative, che seguono direttamente dall'architettura corrente, e in direzioni di ricerca, che richiedono validazione sperimentale indipendente.

Nel \textbf{Capitolo~\ref{chap:conclusions}} vengono tratte le conclusioni della tesi. Per ciascuno dei tre gap identificati nell'introduzione si verifica come l'obiettivo corrispondente sia stato raggiunto, valutandone la portata effettiva e collocando il contributo nel contesto più ampio della sicurezza come processo continuo e verificabile.